
\setcounter{footnote}{26}

\medskip
\noindent\SGAInlineRefTarget{sga3:iii:corollary:0-7-bis}\textbf{Corollary 0.7 bis.}
\textit{Let \(X\) be an \(S\)-scheme such that \(X_{\mathscr J}\) is an
\(S_{\mathscr J}\)-monoid. Then \(L_X\) has the structure of an
\(S\)-monoid, and there is a split exact sequence of \(S\)-monoids:}

\SGAThreeNativeDiagram[0.50\textwidth]
  {native_diagrams/exp3/diagram_024_exp3_p017_corollary07bis_split_sequence.tex}
  {The split exact sequence of \(S\)-monoids in \SGARefLink{sga3:iii:corollary:0-7-bis}{Corollary 0.7 bis}.}
  {fig:exp3-corollary-07bis-split-sequence}
  {Polo--Gille Exp3-14oct24.pdf, component-local p. 17; combined-reader p. 127; printed p. 117}

\textit{The monoid law induced on \(L'_X\) coincides with its abelian-group
structure. In particular, if \(X_{\mathscr J}\) is an
\(S_{\mathscr J}\)-group, then \(L_X\) is an \(S\)-group and is the
semidirect product of \(X^+\) and \(L'_X\).}

\noindent\textit{Proof.}
Since \(X_{\mathscr J}\) is an \(S_{\mathscr J}\)-monoid,
\[
 X^+=\prod_{S_{\mathscr J}/S}X_{\mathscr J}
\]
is an \(S\)-monoid; indeed,
\(X^+(Y)=X_{\mathscr J}(Y_{\mathscr J})\) for every \(Y\to S\).
For every \(S\)-scheme \(Y\), let \(\widetilde Y_{\mathscr J}\) denote
the affine \(Y_{\mathscr J}\)-scheme corresponding to the quasi-coherent
\(\mathscr O_{Y_{\mathscr J}}\)-algebra
\[
 \mathscr O_{Y_{\mathscr J}}\oplus\mathscr J\mathscr O_Y,
\]
that is, the graded algebra associated with the filtration
\(\mathscr O_Y\supset\mathscr J\mathscr O_Y\). Then \(L_X(Y)\) identifies
with \(X_{\mathscr J}(\widetilde Y_{\mathscr J})\), and \(L'_X(Y)\)
identifies with the kernel of the morphism
\[
 p:X_{\mathscr J}(\widetilde Y_{\mathscr J})
   \longrightarrow X_{\mathscr J}(Y_{\mathscr J})
\]
induced by the \emph{zero section}
\(Y_{\mathscr J}\to\widetilde Y_{\mathscr J}\). Equivalently, that
section is induced by the morphism of
\(\mathscr O_{Y_{\mathscr J}}\)-algebras
\[
 \mathscr O_{\widetilde Y_{\mathscr J}}
   \longrightarrow\mathscr O_{Y_{\mathscr J}}
\]
that vanishes on the ideal \(\mathscr J\mathscr O_Y\). Thus, for every
\(Y\to S\), there is a split exact sequence of monoids, functorial in
\(Y\):

\SGAThreeNativeDiagram[0.60\textwidth]
  {native_diagrams/exp3/diagram_025_exp3_p017_corollary07bis_pointwise_split_sequence.tex}
  {The pointwise split exact sequence for \(L_X(Y)\).}
  {fig:exp3-corollary-07bis-pointwise-sequence}
  {Polo--Gille Exp3-14oct24.pdf, component-local p. 17; combined-reader p. 127; printed p. 117}

It remains to show that the monoid law induced on \(L'_X\) coincides
with its abelian-group structure. Let \(\mu\) be the monoid law of
\(L_X\), and let \(e\) be its unit section. We must prove that, for
every \(m,m'\in L'_X(Y)\),
\[
 \mu(m\mathbin{\cdot}e,m'\mathbin{\cdot}e)
   =(m+m')\mathbin{\cdot}e.
\]

This may be seen in either of two ways. First, repeat the proof of
(\SGARefLink{sga3:iii:corollary:0-1-10}{0.1.10} \((\ast)\)), replacing the morphism \(f:X\to W\) there by
\[
 \mu:L_X\times_S L_X\longrightarrow L_X.
\]
Identify \(X^+(Y)=X_{\mathscr J}(Y_{\mathscr J})\) with its image under
\(s\) in \(L_X(Y)=X_{\mathscr J}(\widetilde Y_{\mathscr J})\). For all
\(g,g'\in X_{\mathscr J}(Y_{\mathscr J})\) and
\(m,m'\in L'_X(Y)\), one then obtains
\begin{equation}
 \mu(m\mathbin{\cdot}g,m'\mathbin{\cdot}g')
 =L_\mu^{(g,g')}(m,m')\mathbin{\cdot}\mu(g,g').
 \tag{$\star$}\label{sga3:iii:corollary:0-7-bis:equation:star}
\end{equation}
Here \(L_\mu^{(g,g')}\) denotes the morphism derived from \(\mu\) at the
point \((g,g')\); in other words,
\(\widetilde Y_{\mathscr J}\) lies over \(L_X\times_S L_X\) through
\((g,g')\). In particular,
\[
 \mu(m\mathbin{\cdot}e,m'\mathbin{\cdot}e)
 =L_\mu^{(e,e)}(m,m')\mathbin{\cdot}e.
\]
But
\[
 L_\mu^{(e,e)}(m,m')=L_{\ell_e}(m')+L_{r_e}(m),
\]
where \(\ell_e\), respectively \(r_e\), denotes left, respectively
right, translation by \(e\). Each is the identity map of
\(X_{\mathscr J}\), and therefore
\(L_\mu^{(e,e)}(m,m')=m+m'\).

Alternatively, one may argue as follows (cf. the proof of [DG70],
\S II.4, \SGARefLink{sga3:iii:theorem:3-5}{Theorem~3.5}). By \SGARefLink{sga3:iii:lemma:0-1-11}{Lemma~0.1.11}, formation of \(X^+\) and of
\(L_X\) commutes with products, and hence so does formation of \(L'_X\).
It follows that the morphism
\[
 \mu':L'_X\times L'_X\longrightarrow L'_X
\]
induced by \(\mu\) is a homomorphism for the abelian-group structure.
\SGARefLink{sga3:iii:corollary:3-10}{Lemma~3.10} of Expos\'e~II then shows that \(\mu'\) coincides with the
abelian-group law.

\medskip
\noindent\SGAInlineRefTarget{sga3:iii:paragraph:0-8}\textbf{0.8.}\footnote{Editor's note: the number 0.8 has been
inserted here to mark the return to the original text.}
Now let \(X\) be an \(S\)-scheme such that \(X_{\mathscr J}\) is an
\(S_{\mathscr J}\)-group. Suppose that there is an \(S\)-morphism
\[
 P:X\times_S X\longrightarrow X
\]
whose base change
\[
 P_{\mathscr J}:
 X_{\mathscr J}\times_{S_{\mathscr J}}X_{\mathscr J}
 \longrightarrow X_{\mathscr J}
\]
is the group law of \(X_{\mathscr J}\). An important special case is
that in which \(X\) is an \(S\)-group and \(P\) is its group law. We
obtain a morphism
\[
 L_P:L_X\times_S L_X\simeq L_{X\times_S X}\longrightarrow L_X.
\]
In fact, this morphism does not depend on \(P\), since it can be
computed from the group law \(P_{\mathscr J}\) of \(X_{\mathscr J}\),
as we now show.\footnote{Editor's note: the original has been expanded
in what follows.}

By \SGARefLink{sga3:iii:proposition:0-2:item:ii}{Proposition~0.2(ii)} and (iii), for every \(Y\to S\), every
\(x,x'\in X(Y)\), and every \(m,m'\in L'_X(Y)\),
\[
\begin{aligned}
 P(m\mathbin{\cdot}x,m'\mathbin{\cdot}x')
 &=P\bigl((m,m')\mathbin{\cdot}(x,x')\bigr)\\
 &=L_P^{(x,x')}(m,m')\mathbin{\cdot}\mu(g,g'),
\end{aligned}
\]
where \(g\), respectively \(g'\), is the image of \(x\), respectively
\(x'\), in \(X^+(Y)\). Moreover (cf. the proof of \SGARefLink{sga3:iii:corollary:0-10}{0.10}),
\(L_P^{(x,x')}=L_\mu^{(g,g')}\), and by \SGARefLink{sga3:iii:corollary:0-7-bis}{0.7 bis} \((\star)\) this is the
element of \(L'_X(Y)\) defined by the equality
\[
 L_\mu^{(g,g')}(m,m')\mathbin{\cdot}\mu(g,g')
   =\mu(m\mathbin{\cdot}g,m'\mathbin{\cdot}g')
\]
in \(L_X(Y)\).

Write \(\times\), rather than \(\mu\), for the group law of \(L_X\),
and write \(\operatorname{Ad}\) for the \emph{adjoint action} of \(X^+\)
on \(L'_X\). This action factors through \(X_0\) and is induced by the
adjoint action of \(X_0\) on \(\omega^1_{X_0/S_0}\). Then
\[
\begin{aligned}
 L_\mu^{(g,g')}(m,m')\times g\times g'
 &=m\times g\times m'\times g'\\
 &=\bigl(m\times\operatorname{Ad}(g)(m')\bigr)\times g\times g',
\end{aligned}
\]
and hence
\[
 L_P^{(x,x')}(m,m')=m\times\operatorname{Ad}(g)(m').
\]
We have therefore proved:

\medskip
\noindent\SGAInlineRefTarget{sga3:iii:proposition:0-8}\textbf{Proposition 0.8.}
\textit{Let \(P:X\times_S X\to X\) be an \(S\)-morphism such that
\(P_{\mathscr J}\) equips \(X_{\mathscr J}\) with an
\(S_{\mathscr J}\)-group structure. Let \(\times\) denote the group law
of \(L'_X\), let
\((m,x)\mapsto m\mathbin{\cdot}x\) denote the morphism
\(L'_X\times_S X\to X\) defining the action of \(L'_X\) on \(X\), and
let}
\[
 \operatorname{Ad}:X^+\longrightarrow
 \operatorname{Aut}_{S\text{-}\mathrm{gr}}(L'_X)
\]
\textit{be the adjoint action of \(X^+\) on \(L'_X\), induced by the
adjoint action of \(X_0\) on \(\omega^1_{X_0/S_0}\). Then, for every
\(S'\to S\), every \(x,x'\in X(S')\), and every
\(m,m'\in L'_X(S')\),}
\begin{equation}
 P(m\mathbin{\cdot}x,m'\mathbin{\cdot}x')
 =\bigl(m\times\operatorname{Ad}(p_X(x))(m')\bigr)
    \mathbin{\cdot}P(x,x').
 \tag{0.8.1}\label{sga3:iii:proposition:0-8:equation:0-8-1}
\end{equation}

\textit{If \(X\) is an \(S\)-group, we shall write \(\ast\) for its law,
\(e\) for its unit section, and \(i\) for the \(S\)-morphism defined by}
\[
 i(m)=m\mathbin{\cdot}e
\]
\textit{for every \(S'\to S\) and every \(m\in L'_X(S')\).}

\medskip
\noindent\SGAInlineRefTarget{sga3:iii:corollary:0-9}\textbf{Corollary 0.9.}
\textit{Let \(X\) be an \(S\)-group. Then \(X^+\) is naturally an
\(S\)-group, and \(p_X\) is a morphism of \(S\)-groups. Moreover, the
\(S\)-morphism}
\[
 i:L'_X\longrightarrow X,
 \qquad m\longmapsto m\mathbin{\cdot}e,
\]
\textit{is an isomorphism of \(S\)-groups from \(L'_X\) onto
\(\operatorname{Ker}(p_X)\), and, for every \(S'\to S\), every
\(x'\in X(S')\), and every \(m\in L'_X(S')\),}
\begin{equation}
 m\mathbin{\cdot}x'
 =(m\mathbin{\cdot}e)\ast x'
 =i(m)\ast x'.
 \tag{0.9.1}\label{sga3:iii:corollary:0-9:equation:0-9-1}
\end{equation}

The first two assertions were already proved in \SGARefLink{sga3:iii:remark:0-1-2}{0.1.2}. Since \(X\) is
formally principal homogeneous over \(X^+\) under
\(L_X=L'_X\times_S X^+\), the morphism \(i\) is indeed an isomorphism
of \(S\)-functors from \(L'_X\) onto the kernel of \(p_X\). The fact
that \(i\) is a group morphism and formula (\SGARefLink{sga3:iii:corollary:0-9:equation:0-9-1}{0.9.1}) follow from (\SGARefLink{sga3:iii:proposition:0-8:equation:0-8-1}{0.8.1}),
applied respectively with \(x=x'=e\), and with \(x=e,\ m'=1\).

\medskip
\noindent\SGAInlineRefTarget{sga3:iii:corollary:0-10}\textbf{Corollary 0.10.}
\textit{Let \(X\) be an \(S\)-group. With the preceding notation, for
every \(S'\to S\), every \(x\in X(S')\), and every
\(m'\in L'_X(S')\),}
\begin{equation}
 x\ast i(m')\ast x^{-1}
 =i\bigl(\operatorname{Ad}(p_X(x))(m')\bigr).
 \tag{0.10.1}\label{sga3:iii:corollary:0-10:equation:0-10-1}
\end{equation}
This follows from the equality
\(i(m')\ast x^{-1}=m'\mathbin{\cdot}x^{-1}\) and from (\SGARefLink{sga3:iii:proposition:0-8:equation:0-8-1}{0.8.1}), applied with
\(m=1\) and \(x'=x^{-1}\).

Thus, when \(X\) is an \(S\)-group, we have explicitly determined the
kernel of \(X\to X^+\) and the action of the inner automorphisms of
\(X\) on this kernel. We shall now see that the same can be done for
certain, not necessarily representable, \(S\)-group functors. One case
will be useful to us, namely the functors
\(\underline{\operatorname{Aut}}\) of \SGARefLink{sga3:i:subsection:1:7}{Expos\'e~I, 1.7}. We state it at
once in the next proposition.

